%=============================================================================
% Mass_Paper_Error_Fixes.tex
% Three critical corrections for Fermion_Mass_Paper_Draft.tex
% Date: 2026-03-23
%=============================================================================

\documentclass[11pt]{article}
\usepackage{amsmath,amssymb}

\begin{document}

\section*{Error Fix 1: $G_{11}$ Arithmetic}

\subsection*{Current text (lines 248--258)}

\begin{quote}
\verb|\subsection{The Generation Operator $G$}|

The generation operator $G = \mathrm{diag}(2/3,\,1,\,1)$ acts on
$\mathcal{H}_{\text{gen}} = \mathrm{span}\{|1\rangle,|2\rangle,|3\rangle\}$.

The entry $G_{11} = 2/3$ has its origin in the Cayley graph of $A_5$ with
generating set $S = \{a,a^{-1},b,b^{-1}\}$ where $a=(123)$ and $b=(12345)$.
Of the four generators, exactly two ($a$ and $a^{-1}$) have order~3, giving
\[
G_{11} = \frac{|\{s\in S : \text{ord}(s) = 3\}|}{|S|} = \frac{2}{4} = \frac{2}{3}\,.
\]
\end{quote}

\subsection*{Problem}

$2/4 = 1/2$, not $2/3$. The final value $G_{11} = 2/3$ is correct (it is
derived from the primed microsector trace in Theorem~\ref{thm:G-derived}
later in the paper), but the intermediate arithmetic shown here is wrong.

\subsection*{Replacement text}

\begin{quote}
\verb|\subsection{The Generation Operator $G$}|

The generation operator $G = \mathrm{diag}(2/3,\,1,\,1)$ acts on
$\mathcal{H}_{\text{gen}} = \mathrm{span}\{|1\rangle,|2\rangle,|3\rangle\}$.

The entry $G_{11} = 2/3$ is derived from the primed microsector trace
(Theorem~\ref{thm:G-derived}; see also Ref.~[DFDUnified], Theorem~K.4).
The 9-dimensional isotypic block $\Pi$ carrying the generation structure
decomposes into three rank-3 generation projectors $\{M_0, M_1, M_2\}$.
The primed trace prescription removes the self-generation channel, giving
\begin{equation}
G_{11} = \frac{\mathrm{Tr}(\Pi - M_0)}{\mathrm{Tr}(\Pi)}
       = \frac{9 - 3}{9} = \frac{2}{3}\,.
\end{equation}
By normalization convention, $G_{22} = G_{33} = 1$ (generations 2 and 3
couple at full strength at the Higgs vertex).
\end{quote}


%=============================================================================
\section*{Error Fix 2: $b$-Quark Exponent Exception}

\subsection*{Current text (lines 165--206)}

The paper presents the exponent formula as a boxed theorem:
\[
n_f = \frac{k_f + k_H}{2}
\]
and then in Table~1 writes the $b$-quark entry as
``$n_f = 1 \to 0$'' with the note ``shifted to $n=0$ by QCD dressing.''

The formula is stated as universal (Theorem~2), but the $b$ quark is
an explicit exception: $k_b = 1$, $k_H = +1$ gives $n_b = 1$, yet
the v67 dictionary has $n_b = 0$.

\subsection*{Problem}

If the formula requires a QCD correction for one fermion, it is not
universal as stated. The paper must be honest about this.

\subsection*{Replacement text for the theorem (lines 167--178)}

Replace:
\begin{quote}
\begin{verbatim}
\begin{theorem}[$\alpha$-Exponent from Bundle Degree]
\label{thm:exponent}
The Yukawa coupling for fermion species $f$ has the form
$y_f \propto \alpha^{n_f}$ with
\begin{equation}
\boxed{\quad n_f = \frac{k_f + k_H}{2}\quad}
\label{eq:exponent}
\end{equation}
where $k_f\in\Z$ is the fermion bundle degree on $\CP^2$ and
$k_H = +1$ for $H$-coupling (leptons, down-type quarks) or $k_H = -1$
for $\widetilde{H}$-coupling (up-type quarks).
\end{theorem}
\end{verbatim}
\end{quote}

With:
\begin{quote}
\begin{verbatim}
\begin{theorem}[$\alpha$-Exponent from Bundle Degree]
\label{thm:exponent}
The Yukawa coupling for fermion species $f$ has the form
$y_f \propto \alpha^{n_f}$ with bare exponent
\begin{equation}
\boxed{\quad n_f^{\mathrm{bare}} = \frac{k_f + k_H}{2}\quad}
\label{eq:exponent}
\end{equation}
where $k_f\in\Z$ is the fermion bundle degree on $\CP^2$ and
$k_H = +1$ for $H$-coupling (leptons, down-type quarks) or $k_H = -1$
for $\widetilde{H}$-coupling (up-type quarks).

For eight of the nine charged fermions, the physical exponent equals
the bare value: $n_f = n_f^{\mathrm{bare}}$.  The exception is the
$b$~quark, where strong-interaction dressing shifts the exponent:
\begin{equation}
n_b^{\mathrm{bare}} = \frac{1 + 1}{2} = 1 \;\;\xrightarrow{\text{QCD}}\;\;
n_b^{\mathrm{phys}} = 0\,.
\end{equation}
Physically, the $b$~quark sits at the Higgs vertex ($k_b = 1$) in the
$H$-coupling channel, giving a bare exponent of~1.  However, the
large QCD correction from the down-type operator $Q_d$
(with $Q_d[3,3] = 1/42$ encoding $N_f \cdot b_0 = 42$) absorbs the
full $\alpha^1$ suppression, effectively shifting $n_b$ to~0.  The
canonical (QCD-dressed) exponents used in all mass predictions are
those listed in Table~\ref{tab:bundles}.
\end{theorem}
\end{verbatim}
\end{quote}

\subsection*{Replacement text for Table 1, $b$-quark row (line 198)}

Replace:
\begin{quote}
\verb|$b$ & 1 & $+1$ & 1$\to$0  & At vertex; shifted to $n=0$ by QCD dressing \\|
\end{quote}

With:
\begin{quote}
\verb|$b$ & 1 & $+1$ & 1 (bare) $\to$ 0 (phys.) & Bare $n=1$; QCD dressing|\\
\verb|    &   &      &                           & ($Q_d[3,3]=1/42$) shifts to $n_b^{\rm phys}=0$ \\|
\end{quote}


%=============================================================================
\section*{Error Fix 3: Misleading $\alpha^{-1}$ Statement}

\subsection*{Current text (lines 588--598)}

\begin{quote}
\verb|\subsection{Relation to the $\alpha$ Derivation}|

The fine-structure constant itself is derived separately in the DFD framework
from the Chern--Simons coefficient $k_{\max} = 60$ via
$\alpha^{-1} = k_{\max}\cdot(1 + \mathcal{O}(\alpha))$.  Thus the full chain
is:
\[
\text{CP}^2\text{ topology} \;\xrightarrow{\text{heat kernel}}\; b = 60
\;\xrightarrow{\text{Chern--Simons}}\; \alpha \approx 1/137
\;\xrightarrow{\text{spin}^c + A_5}\; \text{9 masses}\,.
\]
\end{quote}

\subsection*{Problem}

With $k_{\max} = 60$ and $\alpha^{-1} \approx 137$, writing
``$\alpha^{-1} = k_{\max}(1 + \mathcal{O}(\alpha))$'' is misleading.
The correction factor is $137/60 \approx 2.28$, which is not a small-$\alpha$
perturbative correction --- it more than doubles the leading term.

\subsection*{Replacement text}

\begin{quote}
\begin{verbatim}
\subsection{Relation to the $\alpha$ Derivation}

The fine-structure constant itself is derived \emph{ab initio} within the
DFD framework from the Chern--Simons level $k_{\max} = 60$ and the
hypercharge trace $\mathrm{Tr}(Y^2)$ (Ref.~[AbInitio]):
\begin{equation}
\alpha^{-1} = \frac{\pi^{3/2}}{24}\,\mathrm{Tr}(Y^2)\,k_{\max}\,
\frac{k_{\max}+3}{k_{\max}+4}\,
\biggl[1 + \frac{7}{80\bigl((k_{\max}+4)^2 - 1\bigr)}\biggr]
= 137.035\,999\,854\,.
\label{eq:alpha-full}
\end{equation}
Thus the full derivation chain is:
\[
\text{CP}^2\text{ topology} \;\xrightarrow{\text{heat kernel}}\;
k_{\max} = 60 \;\xrightarrow{\;\eqref{eq:alpha-full}\;}\;
\alpha^{-1} = 137.036
\;\xrightarrow{\text{spin}^c + A_5}\; \text{9 masses}\,.
\]
\end{verbatim}
\end{quote}


\end{document}
